\section{Purpose, scope and operating model}\label{sec:scope}
Rocket GNC Monitor is a native desktop ground application for an amateur rocket.
It combines two independent serial links, Digital and Analog USB video, recorded-flight playback,
a visual antenna assembly, telemetry and actuator displays, and a nominal
OpenRocket trajectory. The current transport contract is Zephyrus. The interface
has places for a configurable number of canards and four fin tabs, but the
Zephyrus packet contains four legacy servo drives; it does not provide the new
vehicle's requested and measured canard/tab angles.

The application starts in \ui{LIVE}, on \ui{Control panel}, without opening either
serial port, starting polling or opening a camera. Connect each board explicitly.
A ground-board USB connection, a working rocket radio link and a pointer-board
USB connection are three separate observations. A screen showing telemetry does
not imply that an uplink command was acknowledged by the rocket.

\begin{longtable}{@{}L{.13\linewidth}L{.38\linewidth}L{.42\linewidth}@{}}
\toprule Mode & Source and controls & Persistence and transition behavior \\\midrule\endhead
LIVE & Physical Zephyrus telemetry; ground commands require its board connection. Pointer manual controls use an independent physical or virtual connection. & Logging captures received data. Source changes close physical transports and clear displayed history. Reconnect and restart polling explicitly. \\
DEMO & One complete bundled GS1, GS2 or GS3 receiver log; synthetic video test pattern. Commands are simulated. & No physical serial transport. Can record a rehearsal or export the complete selected dataset. Returning to LIVE restores the prior reference and retains mission settings. \\
REPLAY & Indexed session telemetry and recorded video, initially paused. Physical commands disabled. & Seek, step, change speed and export. Session files are read for replay; a new flight archive can preserve the session and displayed position. \\
\bottomrule
\end{longtable}

The three source modes are not launch states. Rocket states such as Preflight,
Flight and Main come from telemetry and are shown within a mode. Likewise,
\ui{Start Logging} is independent of \ui{Start Polling}: opening a recording does
not by itself create telemetry.

\subsection{Supported machines and offline use}
The current verified build is macOS arm64. The distributable app includes Python,
Qt/PySide6, NumPy, serial support, FFmpeg, the private Java runtime, the team's
OpenRocket engine and bridge, motors, demo recordings and Lucida Grande. A
production machine does not need a separate Python or Java installation.
The macOS app is packaged with a minimum-system declaration of macOS 13; that is
not evidence that every such system has been tested.

Windows has native build entry points and a previously tested x64 package.
The current Mac feature revision has not been rebuilt and qualified on Windows.
Linux-related camera and runtime branches exist, but Linux is not a qualified
release target. Keep the whole Windows package together; its executable alone is
insufficient.

Serial reception, local video input, demonstration, recording/replay, coordinate
conversion and simulation work offline when the package and inputs are present.
Online wind retrieval requires internet access. USB devices may require an
appropriate OS driver, and macOS camera permission must be granted.

\subsection{What the app establishes}
A command marked \emph{sent} has been dispatched on the serial interface; neither
legacy protocol provides the command-level acknowledgment needed to prove the
requested physical action. The antenna graphic follows the last sent target,
not measured motor encoders. Its dimensions are illustrative rather than
surveyed CAD.

The OpenRocket output is a nominal reference. It does not qualify the new
controlled rocket's aerodynamic, actuator or control-loop behavior. Integrated
Zephyrus rotations are displayed as received and are not a validated attitude
solution. Missing measurements are shown as unavailable rather than inferred.

\section{A complete operating workflow}\label{sec:workflow}
\subsection{Before connecting equipment}
\begin{enumerate}
\item Start the app and confirm \ui{LIVE} on the connection panel. Both board
cards should be disconnected.
\item Open \ui{Mission $\rightarrow$ Configure}. Name the mission and enter the
launch origin, or choose \ui{URRG}. Enter the ellipsoid and mean-sea-level
elevations separately, then select the verified legacy height interpretation.
\item In \ui{Mission \& wind}, enter and apply a manual wind profile or retrieve
one for the intended UTC time. Select the rocket's \code{.ork} file and any
additional motor curves. Run the nominal simulation if a reference is needed.
\item Save a planning \code{.rktflight} file so the settings, selected model/motors
and available reference can be reopened together.
\end{enumerate}
\subsection{Receiving and recording}
\begin{enumerate}
\item Refresh the two serial selectors. Assign the ground-station board and
antenna-pointer board to different ports.
\item Connect the ground board and click \ui{Start Polling}. Wait for
\ui{Rocket telemetry live}, then verify packet age and the displayed data.
\item Connect the pointer board if pointing control is needed. Manual pointer
control is independent of the ground board.
\item On \ui{Flight overview}, choose \ui{Find cameras}, select the USB receiver
and press \ui{Start camera}. Confirm that frames are arriving.
\item Click \ui{Start Logging}. It immediately creates a unique session folder.
Check Diagnostics for the recording path and any loss/error indication.
\item For tracking, use \ui{Send to AntPtr} to freeze the receiver GPS, then
\ui{Track live rocket}. Tracking requires both boards, polling and a usable
fresh target. Review the meaning of ZERO and Stop tracking in
Section~\ref{sec:antenna}.
\end{enumerate}
\subsection{After the flight}
Stop tracking and stop logging. Use \ui{Save flight} to create the portable
archive; the save waits for pending recorder/video closure. Reopen the file to
check the recorded session in paused replay. Retain the original session folder
when preserving full capture provenance. Close the app normally to finalize its
workers. A flight saved while logging is still active is a snapshot and can omit
the unfinished video segment.

\section{Global settings}\label{sec:settings}
Open \ui{Rocket GNC Monitor > Settings} on Mac or press Command-comma.
Windows uses \ui{File > Settings} or Ctrl-comma. The window is available in every
workspace and source mode, including disconnected LIVE. Reopening it brings the
existing editor forward.

\begin{longtable}{@{}L{.24\linewidth}L{.69\linewidth}@{}}
\toprule Preference & Behavior \\\midrule\endhead
OpenRocket JAR & Browse to a complete compatible OpenRocket-MIT JAR. \ui{Use bundled} clears the override and shows the bundled release and path. Java remains included with the app. \\
Simulation time limit & 10--1800 s; default 120 s. Engine and time-limit changes apply to the next job; the running job retains its snapshot. \\
Recording folder & Destination for new logging sessions. \ui{Use default} restores the application's sessions directory. Existing files and active recordings retain their locations. \\
Appearance & Dark or Daylight. The \ui{View > Daylight theme} action updates the same persistent preference. \\
\bottomrule
\end{longtable}
\ui{Save} validates and stores changes. \ui{Cancel} discards edits.
\ui{Restore Defaults} changes the form and requires Save to take effect.
A missing or unreadable custom JAR is reported; select the relocated file or use
the bundled engine. The current Mac package includes OpenRocket-MIT v6.2 and
requires no separate OpenRocket installation or runtime download.

Preferences are stored in \code{settings.json} under the per-user application
data directory; the exact path is shown in the window. \code{--data-dir} also
selects a separate settings store. These machine-local paths and preferences
are not exported in mission JSON or portable flight files. Site, wind, rail,
alert thresholds and video offset remain mission-specific. Startup remains
LIVE with hardware closed regardless of the saved preferences.

\section{Navigation and connection control}
The persistent header contains the mission name, source-mode selector,
Start/Stop Polling and Start/Stop Logging. In LIVE and REPLAY the two serial cards
remain visible; controls are enabled according to mode and board state. The
status area provides local wall-clock time, source/status messages and alerts.
The \ui{View} menu selects the daylight theme; this changes contrast, not data.

\begin{longtable}{@{}L{.23\linewidth}L{.70\linewidth}@{}}
\toprule Workspace & Function \\\midrule\endhead
Control panel & Ground USB, rocket radio and pointer USB status; packet age, RSSI, valid/rejected counts and readiness. \\
Flight overview & Main metrics, camera/file video, projected trajectory, vertical profile and attitude instrument. \\
Antenna pointer & Orbitable assembly, manual azimuth/elevation, direction commands, ZERO, ground-GPS freeze and tracking. \\
GNC \& actuators & Three angular-rate and integrated-rotation traces, canard/tab channel slots and legacy servo drives. \\
Mission \& wind & Mission configuration, launch location, wind layers, weather, rocket/motor selection and simulation. \\
Sessions \& replay & Portable flight files, session-folder replay, legacy CSV import, sample export, seeking and event timeline. \\
Diagnostics & Decoded sample JSON, packet counters, recording state, events and alert acknowledgment. \\
Rocket controls & Four subviews: Telemetry, Rocket commands, Power, Recovery \& pyros. \\
\bottomrule
\end{longtable}

\subsection{Rocket radio status}
\begin{longtable}{@{}L{.20\linewidth}L{.38\linewidth}L{.35\linewidth}@{}}
\toprule Status & Meaning & Normal next action \\\midrule\endhead
Disconnected & Ground USB transport is not connected. & Select and connect the telemetry board. \\
Paused & Ground USB is open but telemetry polling is stopped. & Start polling to receive; uplinks remain available. \\
Waiting & Polling is active but no fresh valid packet has established reception. & Check transmitter, receiver, antenna and serial source. \\
Receiving & A checksum-valid LIVE packet arrived within the configured freshness interval. & Monitor packet age and loss indicators. \\
Stale & Previously received valid telemetry has aged beyond that interval. & Restore reception; explicitly restart tracking if it stopped. \\
\bottomrule
\end{longtable}
The default freshness interval is 2 s. Reconnecting the ground board clears prior
reception evidence. A connected USB device alone never proves a rocket radio
link. The trailer carrying ground-receiver GPS has no independent checksum.

\subsection{Serial behavior and command availability}
Both boards use 115200 baud, eight data bits, no parity, one stop bit and no
flow-control override. DTR/RTS follow the original pyserial defaults.
Ground reads use a 1 s timeout; pointer reads use 0.1 s. The selectors reject
assigning the same serial device to both roles, including the usual macOS
\code{/dev/tty.} versus \code{/dev/cu.} alias. A connected or connecting
board's port is removed from the other selector and restored when disconnected.

Stopping polling leaves the port open and leaves rocket uplinks available.
Disconnecting a board closes its transport and clears relevant pending state.
There is no automatic reconnection, retransmission or tracking restart. A
partially written packet is treated as uncertain, not retried.

\section{Flight overview, coordinates and time}\label{sec:flightview}
\subsection{Metrics and plots}
The six main metrics are altitude in metres, velocity in metres per second,
aligned flight time in seconds, battery in volts, receiver RSSI in dBm and sample
age in seconds. Missing values remain blank/unavailable. Battery is the sum of
the three decoded cell voltages; velocity is the legacy reported integrated
velocity rather than a recomputed derivative of the plotted altitude.

The trajectory view offers \ui{Plan: East/North}, \ui{Side: East/Up} and a fixed
projected \ui{Perspective}. The actual trace uses accepted local positions. The
reference is dashed, with a diamond at the simulated position for the same
aligned time. The residual label is the Euclidean three-dimensional distance
between actual and reference positions at that time. It is not a closest-point
distance or an estimate of navigation error.

A reference marker requires aligned time within the reference interval; the
app does not extrapolate before or after the simulation. A trace or last point
can remain visible when new packets are absent, so always read sample age.
The vertical profile overlays reported altitude with the reference's
launch-relative up coordinate. Their comparison depends on using the correct
altitude convention.

\subsection{Altitude conventions}
For LIVE geographic-to-local conversion, configure an origin and a valid 3D
rocket GPS fix. Choose one verified interpretation:
\begin{longtable}{@{}L{.27\linewidth}L{.66\linewidth}@{}}
\toprule Mission choice & Conversion used for the rocket's absolute height \\\midrule\endhead
Unknown & Geographic trajectory position remains unavailable. Telemetry values still display. \\
GPS absolute ellipsoid & Use the reported GPS altitude directly. \\
GPS zeroed at launch origin & Add mission ellipsoid altitude to the reported GPS-relative altitude. \\
Filtered barometer relative to launch & Add mission ellipsoid altitude to the reported filtered barometric altitude. \\
\bottomrule
\end{longtable}
Mean-sea-level elevation is a separate simulation/weather input; it is not
substituted for ellipsoid height in the geographic transform. The legacy antenna
tracking convention is separate again (Section~\ref{sec:antenna}).

\subsection{Clock meanings}
There are three primary clocks. Device time is flight-computer uptime in
seconds after rollover extension. UTC is the packet's host reception time.
Monotonic host time measures freshness and session elapsed time. Flight time is
\[
 t_{\mathrm{flight}}=t_{\mathrm{device}}-t_0.
\]
In LIVE, observing a received Preflight-to-Flight transition establishes
\(t_0\). Connecting mid-flight does not establish liftoff time. A detected device
reboot clears time alignment and displayed history and stops tracking. The
Sessions action \ui{Set displayed sample to flight \(t=0\)} provides explicit
manual alignment.

The attitude display uses the legacy integrated rotations. In GNC, axes are
labeled X/roll, Y/pitch and Z/yaw. The attitude graphic and traces describe
the available legacy data; they do not imply a quaternion navigation estimator.

\section{Digital and Analog video}
Flight overview has independent \ui{Digital} and \ui{Analog} panels. Both receive
video from USB cameras; the labels describe the rocket-to-ground radio links.
In each panel choose \ui{Find}, select a listed USB camera/receiver and press
\ui{Start}. Each panel has its own \ui{Stop} and \ui{File\ldots} controls.
A selected camera is removed from the other dropdown; active and closing
captures also reserve their camera. Select the empty choice to release an
unused selection. Stopping or changing one stream leaves the other running. The app uses AVFoundation on Mac, DirectShow on Windows and
the existing Video4Linux branch on Linux. A local video file can also be selected.
Network-stream URLs are not a dedicated operator input in this version.

The displayed video is fitted into a \(960\times540\) RGB frame with its aspect
ratio preserved. Each stream has one FFmpeg input feeding its display and recording. Camera and
demo sources are compressed for recording; file inputs use stream-copy for the
recorded output. Audio is not captured. Files are stored as approximately
30-second Matroska segments with a CSV time index.

Each video status reports frame age and errors; after two seconds without a
recent frame it adds a stale/frozen indication. Stop closes capture but retains
the final displayed image with an explicit stopped label. A retained image is
not proof of a live stream.

Starting or stopping logging reopens both configured video pipelines so their
recording destinations change. A camera reconnect during a recording uses a
new segment directory rather than overwriting earlier footage. Video lifecycle
work occurs off the GUI thread, but a brief stream interruption is possible.

Replay uses host session timing plus the mission's video offset. A positive
offset delays video relative to telemetry; a negative offset advances it.
Each stream follows its own recorded starts, segments and gaps. A paused seek
decodes one frame per stream. Flight files include both feeds; old single-video
sessions appear in Digital. A panel's Stop pauses that replay feed until a seek
or playback restart. Compact displays can scroll to the lower flight plots.
Camera exposure timestamps are not
synchronized with onboard telemetry in this implementation.

\section{Antenna assembly and pointing}\label{sec:antenna}
\subsection{Viewing the physical model}
The drawing follows the supplied four-leg stand, azimuth turntable, elevation
supports and common antenna beam. The grid reflector, long Yagi and enclosed
Avenger XR18 share the elevation carrier; mounting centers are attached to the
beam concentric with the support pivots. The dish feed arm uses the dish's grid
material and lies on the dashed boresight. The Avenger enclosure represents the
black unit identified in the supplied photograph; no product photograph is
embedded in the renderer.

Drag to orbit the viewing camera, scroll to zoom, and double-click to reset the
perspective. Perspective, Front, Side and Top presets are available. Viewing
works with serial controls locked and never sends a command. All bearings use
the corrected chirality: \(0^\circ\) north, \(90^\circ\) east, \(180^\circ\) south,
\(270^\circ\) west.

\begin{figure}[htbp]\centering
\begin{tikzpicture}
\node[box,text width=35mm] (stand) at (0,0) {Four-leg stand\\fixed world frame};
\node[box,text width=35mm] (turn) at (5,0) {Turntable and supports\\azimuth rotation};
\node[accentbox,text width=35mm] (beam) at (10,0) {Common pivot beam\\azimuth and elevation};
\node[smallbox] (dish) at (5.7,-2) {Grid dish\\and aligned feed};
\node[smallbox] (yagi) at (10,-2) {Long Yagi\\beam attachment};
\node[smallbox] (helix) at (10,-3.8) {Avenger XR18\\beam attachment};
\draw[flow] (stand)--(turn);
\draw[flow] (turn)--(beam);
\draw[flow] (beam.south west)--(dish.north);
\draw[flow] (beam)--(yagi);
\draw[flow] (beam.east) -- ++(.6,0) |- (helix.east);
\node[captionnode,text width=42mm] at (.3,-2.8) {Orbit and zoom change the camera only. They do not change commanded azimuth or elevation.};
\end{tikzpicture}
\caption{Transform hierarchy of the procedural antenna model. Both outer antennas and the reflector belong to the same moving carrier; this diagram describes attachments, not dimensions or a photographic viewing direction.}
\label{fig:mount}
\end{figure}


\subsection{Manual controls and the meaning of ZERO}
Enter finite azimuth/elevation values and press \ui{Send}. Editing a text field
alone does not move the drawing or transmit. The graphic updates on successful
serial dispatch. The legacy manual path does not add mission calibration,
software envelope or cable-wrap prerequisites; stored mission envelope fields
are not enforcement for these buttons.

UP, DOWN, LEFT and RIGHT send the original firmware direction opcodes in
5-degree increments. These are direction packets, not host-generated absolute
targets. Before a known target exists, the app cannot show a reliable absolute
pose after a jog.

\ui{ZERO} sends the firmware's count/target reset at the current physical
position. It does not seek a home switch or mechanically return the mount to
north/horizontal. The drawing's last-sent label is a command estimate; measured
angles and board acknowledgment are unavailable.

\subsection{Ground-GPS freeze and automatic tracking}
With both boards connected and telemetry polling, press \ui{Send to AntPtr}.
The receiver latitude and longitude are frozen at the original five-decimal
rounding; its displayed altitude/fix are retained with that snapshot. The legacy
calculation uses ground altitude zero and rocket filtered barometric altitude.
Mission launch origin, ellipsoid/MSL settings and stored pointer offsets do not
replace this convention.

\ui{Track live rocket} computes targets at up to 5 Hz. It requires fresh LIVE
rocket data, usable rocket coordinates and a frozen ground position. The
geographic pointing helper rejects targets within 10 m and nearly overhead
targets with under 1 m horizontal separation. Stale data, reboot or connection
loss stops host tracking; resume explicitly after correcting the condition.
\ui{Unfreeze GPS} releases the stored origin and stops tracking.

\statusnote{Stop tracking}{This stops new targets and removes an unsent queued
target. A movement already commanded can continue. The old firmware has no
motion-stop opcode, watchdog, measured-pose report or mechanical-homing contract.}

\subsection{Virtual pointer connection and rehearsal}
The pointer selector also offers \ui{Virtual antenna pointer} in LIVE.
Alternatively, click \ui{Connect virtual pointer} on the antenna page.
This creates an in-process simulator and needs no ground station or physical
serial port. The connection, pose and banner explicitly identify simulated
motion; rocket controls remain locked until the real ground station connects.

Configure the mount under \ui{Mission > Configure > Antenna pointer}: WGS84
latitude, longitude and ellipsoid altitude, plus \ui{Antenna location established}.
Mission JSON and portable flight files retain the position. The physical-board
tracking path continues using the frozen receiver GPS from \ui{Send to AntPtr}.

Run OpenRocket or load a georeferenced, launch-relative ENU reference. Then use
\ui{Follow trajectory} on the antenna page. Pause/resume, Rewind, a time slider
and 0.25--4 times speed control an independent rehearsal clock. The gold target
and purple antenna marker in Flight overview identify the simulated geometry;
telemetry samples, flight counters and actual tracks are unchanged.

Send, the four direction controls, ZERO and their original keyboard shortcuts
actuate the simulated mount. Manual movement pauses playback. The displayed
pose approaches its target at illustrative limits of 90 degrees/s in azimuth
and 60 degrees/s in elevation. \ui{Hold virtual pointer} freezes the simulated
motion and clock. These rates are visualization assumptions, not measured mount
dynamics. Disconnect or a source-mode change closes the simulator. Editing the
antenna location or replacing the reference stops rehearsal until an explicit
restart. Connection and rehearsal state are not restored from a saved flight.
\begin{figure}[htbp]\centering
\begin{tikzpicture}
\node[box,text width=36mm] (reference) at (0,0) {OpenRocket reference\\ENU path + launch origin};
\node[accentbox,text width=36mm] (frame) at (5.15,0) {Antenna-local direction\\ENU via ECEF};
\node[box,text width=36mm] (pose) at (10.3,0) {Virtual pointer\\command + slew model};
\node[box,text width=36mm] (mission) at (0,-2) {Mission profile\\antenna lat/lon/height};
\node[box,text width=36mm] (clock) at (5.15,-2) {Independent clock\\play, pause, seek, speed};
\node[box,text width=36mm] (display) at (10.3,-2) {Simulated pose + markers\\telemetry unchanged};
\draw[flow] (reference) -- (frame);
\draw[flow] (frame) -- (pose);
\draw[flow] (mission.east) -- (frame.south west);
\draw[flow] (clock) -- (frame);
\draw[flow] (pose) -- (display);
\end{tikzpicture}
\caption{Virtual rehearsal uses the saved antenna position and its own clock. No physical serial transport participates in this path.}
\end{figure}


\section{Rocket commands, power and recovery}\label{sec:rocket}
\subsection{Telemetry subview}
The legacy Telemetry subview retains the original values: state, RSSIs, packet
age/number, filtered and maximum altitude, velocity, integrated rotations,
angle from vertical, temperature, GPS fix/position/heights/precision/satellites,
four servo drive/degree readouts and three cell voltages. Cells above 3.7 V use
the normal accent, those above 3.5 V the intermediate color, and the remainder
the low-value color. These cell cues are separate from the configurable
total-battery alert.

\subsection{Rocket commands subview}
\begin{longtable}{@{}L{.28\linewidth}L{.65\linewidth}@{}}
\toprule Control & Implemented behavior \\\midrule\endhead
Advance State & Requests current decoded state code plus one; retains the original confirmation dialog. The UI does not implement an independent flight-state machine. \\
Zero P/Y/Roll; Zero Alt; Zero Velo; Zero Servos & Sends the corresponding original opcode. These do not rewrite recorded historical samples. \\
PD Activate & Sends the legacy activation command. No new controller tuning interface is implemented. \\
Roll control angle; Airbrakes angle & Sends the entered finite angle as the original float32 payload. \\
VTX 1, 3, 5 or 8 W & Sends the legacy power-choice index. A selected button represents the request, not a radio acknowledgment. \\
\bottomrule
\end{longtable}
Ground commands require the ground board in LIVE even when polling is paused.
They are simulated in DEMO and disabled in REPLAY. The app does not require
fresh radio evidence before every uplink; the separate reception indicator must
be interpreted accordingly.

\subsection{Power subview}
Six converter requests and voltage/current readbacks are shown for 3, 3.3, 5,
7.4, 8.4 and 28 V. Changing an editable request sends the full six-bit mask.
The 3.3 V request stays on and is not editable, matching the old UI.
Total current and temperature are readouts, not switchable rails.
Readback state is derived from voltage proximity to the nominal rail, not from
the state of its checkbox.

The BMS table displays Enabled and Triggered rows for SCD, OCD2, OCD1, OCC,
COV, CUV and the two reserved bits. It is read-only in the original operator
workflow. A BMS command exists at protocol level but is not an additional
operator toggle.

\subsection{Recovery and pyro subview}
Six channels, numbered 0--5, have selections, continuity, armed/fired state and
resistance readouts. Select at least one channel before using ARM or FIRE.
Both retain the original confirmation behavior and encode the chosen channel
mask. The emergency actions PISTON, BP WELLS, TNDR DSNDR and DEPLOY ALL each
retain their own confirmation.

There is no added Disarm button in this subview, although the packet encoder
supports the legacy disarm opcode. Continuity and armed/fired displays come from
received packets; sending a request alone does not update them as confirmed
hardware state. Consult the team's operational procedures for physical recovery
system work; this report documents the existing UI rather than a firing sequence.

\section{Mission configuration and launch sites}\label{sec:mission}
\ui{Mission $\rightarrow$ Configure} opens a scrollable configuration dialog.
Saving validates the entries and applies them; Cancel leaves the current mission
unchanged. Stop logging before editing or loading another mission so a recording
is not redefined in progress. Fields include the mission name, origin-established
flag, location, two altitude datums, legacy height interpretation, canard count,
freshness/battery thresholds, launch rail, saved simulation index, random seed and
video offset. Appendix~\ref{app:mission} lists every stored field, including
legacy fields that are not active UI controls.

\subsection{Latitude/longitude, MGRS and Plus Codes}
All coordinate conversions are offline:
\begin{itemize}
\item \ui{Latitude / longitude}: decimal degrees, latitude \(-90\) to \(90\),
longitude \(-180\) to \(180\).
\item \ui{MGRS}: compact or spaced UTM/UPS grid reference. The preview identifies
the supplied grid resolution. Conversion uses the grid reference point, not a
newly estimated survey position.
\item \ui{Plus Code}: a full code resolves to its cell center. A short code
requires explicit nearby reference latitude and longitude; locality names are
not geocoded.
\end{itemize}
The preview shows the resolved canonical WGS84 coordinates before Save.
Changing only the display format preserves canonical coordinates and avoids
round-trip drift. A saved mission keeps the chosen format/code as entry
provenance alongside latitude and longitude.

\subsection{URRG preset}
Select \ui{URRG} from the launch-site preset menu. It sets the exact supplied
MGRS string \code{18TUN2061530290}, switches the entry format to MGRS and marks the
origin established. Its resolved coordinates are approximately
\(42.7041678^\circ\) latitude and \(-77.1902950^\circ\) longitude.

The preset supplies coordinates only. It does not infer launch elevation,
antenna position, wind or rail settings. Existing elevation values remain for
the operator to review. Manual coordinate edits clear the preset name to Custom;
changing the display format alone keeps URRG. Mission JSON and flight archives
preserve the site selection.

\subsection{Mission JSON versus portable flight files}
\ui{Save mission} writes configuration JSON. \ui{Open mission} restores it,
clears history/track and deselects the old reference. The JSON contains file
paths for a selected rocket and custom motors; it does not embed those files.
Loading mission JSON is not the same operation as opening a flight archive and
does not itself perform the archive workflow's full source disconnection.

Use \ui{Save flight} for a portable package containing the referenced assets
and available data. Keep the distinction in mind when moving work to another
machine.

\section{Wind and nominal OpenRocket simulation}\label{sec:simulation}
\subsection{Manual and imported wind}
The wind table uses height above launch in metres, speed in metres per second
and meteorological direction \emph{from} true north in degrees. Add or remove
layers and press \ui{Apply wind}. Editing a table cell alone does not update the
active mission profile.

Heights must be nonnegative and strictly increasing; there must be 1--200 layers.
Speeds must not exceed 150 m/s. A zero-speed direction has no physical effect.
Internally each row becomes east and north velocity components:
\[
 v_E=-s\sin\theta,\qquad v_N=-s\cos\theta .
\]
\ui{Import wind} and \ui{Export wind} use JSON with \code{layers} and a
\code{source} description. The included example shows the exact format.
Outside the supplied altitude coverage, the nearest wind layer is held.

\subsection{Online weather}
Enter an ISO 8601 time with a timezone, such as
\code{2026-09-21T16:00:00Z}, and the launch elevation above mean sea level.
\ui{Retrieve wind} requests an Open-Meteo pressure-level profile at the mission
latitude/longitude. The current implementation selects forecast or historical
forecast based on age, rejects requests more than 15 days ahead, and requires
an available provider hour within one hour of the selected time. These are app
rules, not a guarantee that the provider has every date/location.

Pressure levels are 1000, 925, 850, 700, 500, 300, 200 and 100 hPa. Missing levels
and below-launch levels are skipped; remaining layers are sorted.
Geopotential height is converted to geometric height, then the launch MSL
elevation is subtracted. The provider response, request, endpoint, retrieval time
and attribution remain in the mission for offline reuse. Review the resulting
coverage. If the launch coordinates change during retrieval, the UI rejects
applying that result to the new site.

\subsection{Running a reference}
\begin{enumerate}
\item Configure the launch origin and altitude conventions.
\item Choose the team's \code{.ork} file. It must contain a saved simulation with
a valid active motor configuration.
\item Select the saved simulation index in mission settings; index 0 means the
first saved simulation.
\item The Zephyrus MITRT N8406 curve is bundled automatically. For other custom
motors, use \ui{Select custom motor files} to provide matching \code{.eng} or
\code{.rse} curves. Cancel keeps the existing selection.
\item Review rail length, tilt from vertical, true-north heading, wind and seed.
Press \ui{Run nominal simulation}. \ui{Cancel simulation} terminates the worker.
\end{enumerate}
A fresh simulation uses the saved vehicle/motor configuration but explicit
app-supplied site, rail and wind settings. The bridge disables saved extensions
and the fork's global inertia override, uses the ISA atmosphere, zero wind
turbulence standard deviation and a fixed seed, and returns the first output
branch. The job also records branch count and warnings.

A missing thrust curve is named before simulation; an unpowered configuration
is rejected. The supplied Zephyrus test model previously failed because its
custom N8406 curve was absent. The bundled curve now matches its required motor
digest. The app copies model/motor inputs into each job and records hashes so
the run can be traced independently of original input paths.

\subsection{Importing a reference and interpreting results}
\ui{Import reference} on Flight overview accepts a CSV and a same-basename JSON
manifest. The CSV columns are \code{time_s,east_m,north_m,up_m}.
The manifest declares version 1, ENU, \code{m,s}, and
\code{origin:[latitude,longitude,ellipsoid_altitude]}; up is launch-relative.
A LIVE comparison requires a configured mission origin matching that reference
within \(10^{-6}\) degrees and 0.1 m altitude. Synthetic references are rejected
for LIVE comparison. Stop logging before replacing the recording's reference.

The saved trajectory is not a flight-control command, and the app does not
run an online controlled-vehicle simulation driven by each received packet.
For GNC assessment, compare the nominal baseline with actual data while retaining
the solver and unavailable-feedback limitations.

\section{Recorded Zephyrus demonstration}
Choose DEMO explicitly. GS2 is the default; GS1 and GS3 are alternative receivers
of the same launch, not consecutive segments. They are never concatenated or
deduplicated. Each plays according to its own receive-time deltas, keeping gaps,
original UTC, packet numbers and device time.

\Needspace{14\baselineskip}
\begin{longtable}{@{}L{.16\linewidth}L{.16\linewidth}L{.18\linewidth}L{.41\linewidth}@{}}
\toprule Receiver & Samples & Duration & Use and limitations \\\midrule\endhead
GS1 & 4,366 & 274.056 s & Original receiver record, with GPS outliers. \\
GS2 & 5,464 & 330.234 s & Default record; most samples and fewer major GPS corruptions. \\
GS3 & 4,939 & 332.725 s & Contains invalid ground GPS and anomalous rocket GPS/state readings; automatic simulated tracking is unavailable. \\
\bottomrule
\end{longtable}
The initial cue is five receive-time seconds before the first FLIGHT packet.
Use Play/Pause, the timeline, 0.25/0.5/1/2/4 times speed, \ui{Launch -5s} or
\ui{Full start}. Stop logging before seeking or changing the receiver log.
Pause and speed changes are allowed during a rehearsal recording.
At the end the final sample is retained; restarting returns to the launch cue.

The demo track uses horizontal GPS offsets from the first FLIGHT position plus
the reported barometric height. It is a local observed track, not a surveyed
ellipsoid/MSL trajectory. Missing 3D fixes and positions over 20 km horizontally
from its origin are omitted from the plot but kept in raw decoded values.
No interpolated packets, touchdown, actuator feedback or corrected state
sequence are invented. The records end in Main descent.

Both videos are labeled \ui{TEST PATTERN} (Digital motion pattern and Analog bars); no actual test-launch video was supplied.
The original 43 CSV fields remain in each decoded sample's provenance. Raw radio
frames were not supplied with these CSVs, so their radio checksums cannot be
revalidated. The app verifies the bundled CSV byte hashes on loading.

\section{Logging, sessions and portable flight files}\label{sec:persistence}
\subsection{Starting and stopping logging}
Start Logging creates a uniquely named session under the application data
directory's \code{sessions} folder. It records accepted samples, events,
raw serial bytes, both legacy CSVs and any recorded video. Logging runs
independently of UI redraws. A bounded recorder queue reports dropped records
and a fault if it overflows; loss is not silently labeled complete.

Stop Logging detaches the active recorder and finishes outstanding video and
file work in the background. The event timeline reports completion or error.
Closing the app also finalizes its workers. A crash can lose uncommitted data and
the active video segment; committed session data may still be replayable.

\subsection{Session-folder replay}
Use \ui{Open recorded session} to select a session folder. Playback begins paused.
Controls provide seek, a 0.1 s step, 0.25/0.5/1/2/4 times speed and manual time
alignment. Seeking reconstructs earlier events and up to 120 s of sample
history instead of retaining future pointer/time state. The complete session
remains indexed on disk even though visible plots are bounded.

\ui{Export samples to CSV} writes the richer normalized export, including elapsed
time and the complete sample JSON. This is distinct from the two original-format
CSV files written automatically by the recorder.

\ui{Import Zephyrus CSV} accepts the audited legacy layout, creates a converted
session in a chosen parent directory, retains the original source file and
opens it for replay. It preserves available extra columns; unsupported layouts
are rejected. Importing decoded CSV cannot recreate raw packet quality or radio
bytes.

\subsection{Save flight and Open flight}
Use \ui{File $\rightarrow$ Save flight} (Command-S on Mac) and
\ui{File $\rightarrow$ Open flight} (Command-O). Both are also available in Mission
and Sessions. A \code{.rktflight} archive contains the mission, embedded selected
rocket/custom-motor files, current normalized simulation reference, and available
session data. It is readable after the original model/motor paths disappear.

\begin{longtable}{@{}L{.28\linewidth}L{.65\linewidth}@{}}
\toprule State when saving & What is included \\\midrule\endhead
Logging active & Consistent committed telemetry/raw/CSV snapshot and finalized video segments. Logging continues. The current unfinished video segment is omitted and the session remains marked as a recording snapshot. \\
Logging just stopped & Last recording in the current source mode; save waits for pending recorder and video finalization. \\
REPLAY & Opened session and displayed replay position. \\
DEMO without a selected recorded session & Complete chosen GS dataset and current demo position; no invented camera footage or raw radio bytes. \\
LIVE without a recorded session & Only the displayed telemetry buffer, up to 6,000 samples, explicitly identified as a buffer export. \\
No available telemetry & Planning flight: mission, model/motors and any current simulation reference. \\
\bottomrule
\end{longtable}
Save again after stopping logging to include the final samples and video segment.
A flight archive carries normalized simulation output and metadata, but is not an
automatic copy of every file in the original Java job directory; keep that
directory for its complete worker log and inputs.

Opening requires logging to be stopped and a running simulation to finish or
be cancelled. The archive is validated before switching the current flight.
A session-bearing flight opens in paused REPLAY with both transports closed.
A planning flight opens in disconnected LIVE. Saved pointer calibration is
cleared; no hardware automatically reconnects.

Saving writes a temporary archive and atomically replaces the destination after
success; failure leaves an existing destination intact. Loading validates paths,
sizes, hashes and supported schemas. A damaged file reports an error instead of
silently supplying partial data. Extracted assets are retained under the app's
\code{flights} cache. The file format is versioned but not encrypted or
cryptographically authenticated.

\subsection{Storage and recovery}
The default root is Qt's per-user \code{AppLocalDataLocation}, under organization
\code{MITRocketTeam} and application \code{RocketGNCMonitor}. It normally lies
under macOS Application Support or Windows Local AppData. Use the paths displayed
by the app rather than assuming a fixed home directory. Developers can override
it with \code{--data-dir}.

A session contains \code{manifest.json}, \code{session.sqlite},
\code{raw.bin}, \code{telemetry.csv}, \code{telemetry_badpackets.csv}, and optional
reference, imported-source CSV and video directories.
Do not delete a crashed session's SQLite \code{-wal} file before recovery:
committed records can reside there. Incomplete sessions remain labeled
incomplete; a usable database is not a guarantee of complete footage or raw bytes.

\section{Alerts, diagnostics and troubleshooting}
Diagnostics shows packet counts, the current decoded sample, event messages and
recording information. Valid/rejected frames, discarded bytes and inferred
sequence gaps describe different conditions. Rejected payloads are isolated
from live sample displays and written to the bad-packet CSV when logging.

The LIVE alert manager tracks telemetry stale/absent, low total battery,
recording fault/loss and independently identified Digital/Analog video ended/stale. Battery must remain below threshold
for 1 s; clearing uses 0.3 V hysteresis. Video alerting has a 1 s dwell after an
error or a frame age over 2 s. Acknowledgment records an event and marks the alert;
it does not clear the underlying condition. The demo can show low original
battery values in readouts without those being a new live-hardware alarm.

\begin{longtable}{@{}L{.30\linewidth}L{.63\linewidth}@{}}
\toprule Symptom & Checks or explanation \\\midrule\endhead
Ground connected, no data & Start Polling. Check that the selected port is the telemetry board and that the rocket/receiver supplies valid frames. USB connection is not radio evidence. \\
Controls remain locked & Confirm mode LIVE and the relevant board connection. Pointer manual controls need the pointer board; rocket commands/logging need the ground board. \\
Camera absent or no image & Find cameras again, confirm the receiver enumerates as a camera, grant permission and inspect the video error. Another process may own the device. \\
Image visible but no motion & Read frame age and stopped/stale status; the last frame can be retained. DEMO uses a test pattern. \\
No actual trajectory & Establish launch origin, select the verified height convention and obtain valid 3D GPS. DEMO filters plot-only GPS outliers. \\
No same-time reference marker & Establish flight time alignment and ensure the displayed time falls within the reference interval. \\
Missing motor / unusable simulation & Select a saved powered configuration and matching custom curve. The N8406 is automatic. Review the named error and job's \code{worker.log}. \\
Tracking unavailable or stopped & Check both links, polling, frozen ground GPS, target freshness and valid coordinates. Restart deliberately after fixing the problem. \\
Graphic disagrees with mount & It is a last-sent command model, not feedback. ZERO resets counts at the physical position. Review firmware sign/mechanical reference separately. \\
Flight save has less video & Active snapshots omit the unfinished segment. Stop logging, wait for completion and save again. \\
Only a short live flight was saved & Logging was not active; the fallback saves only the bounded display buffer. Start Logging before acquisition for a full capture. \\
Flight file rejected & Inspect the error for version, checksum, missing asset, path, database or space problems. Keep the original and the session source; do not bypass validation. \\
Recording incomplete/loss & Inspect recorder error, dropped count and disk availability. Preserve SQLite/WAL and finalized video; neither raw nor CSV can recreate data never received or committed. \\
\bottomrule
\end{longtable}

\section{Keyboard shortcuts and operational limits}\label{sec:shortcuts}
Qt preserves the original shortcut strings. On Mac, Ctrl in those strings maps
to Command, and Alt maps to Option. The following actions work across workspaces;
their enabled state still follows the appropriate connection/mode.

\begin{longtable}{@{}L{.39\linewidth}L{.27\linewidth}L{.27\linewidth}@{}}
\toprule Action & Mac & Windows / Qt notation \\\midrule\endhead
Settings & Command-comma & Ctrl+comma \\
Open flight & Command-O & Ctrl+O \\
Save flight & Command-S & Ctrl+S \\
Refresh serial ports & Command-R & Ctrl+R \\
Ground connection toggle & Command-Option-C & Ctrl+Alt+C \\
Pointer connection toggle & Shift-Command-Option-C & Shift+Ctrl+Alt+C \\
Polling toggle & Command-Return & Ctrl+Return \\
Logging toggle & Command-L & Ctrl+L \\
Pointer UP / DOWN & Command-Up / Down & Ctrl+Up / Down \\
Pointer RIGHT / LEFT & Command-Right / Left & Ctrl+Right / Left \\
Pointer ZERO & Command-0 & Ctrl+0 \\
\bottomrule
\end{longtable}

Field acceptance still requires the actual boards, USB video receiver and mount,
representative data rates, reconnect/power-loss behavior and a launch-duration
soak. Three-axis control feedback, a qualified controlled-vehicle model, measured
pointer pose and firmware-level stop/acknowledgment are outside the presently
verified contract. The current development Mac distribution is ad-hoc signed;
production notarization and clean-machine acceptance remain separate release work.
